\documentclass[10pt]{article}
\usepackage[margin=20mm]{geometry}
\usepackage{amsmath,amssymb,mathtools,fontspec,unicode-math,xurl,hyperref}
\IfFontExistsTF{DejaVu Serif}{\setmainfont{DejaVu Serif}}{\setmainfont{Latin Modern Roman}}
\setmathfont{latinmodern-math.otf}
\setlength{\parindent}{0pt}\setlength{\parskip}{6pt}\sloppy
\title{Universality, Tangential Fibres
and Reductive Coefficients}
\author{AI-assisted technical review}
\date{}
\begin{document}
\maketitle
Abstract: The exact author manuscript arXiv:2109.09301v3 is examined. Four Major proof defects and five Minor corrections are established. The Major defects concern ramified-stalk equivariance, full K3 cohomology, the missing principal-polarization reduction, and the tensor-power step supplying the Artin base case. These defects do not refute the final universality theorem. Explicit replacements are supplied, including a noncircular construction of Artin representations. The remaining checks address proalgebraic completions, tensor powers, equal-characteristic cohomology, polarization, semisimplification, and tangential base points. The scientific decision remains pending.\par
\section*{Source identity and assessment}
Frozen identifier: M25-0258. Author: Alexander Petrov. Article: Universality of the Galois action on the fundamental group of \(\mathbb P^1\setminus\{0,1,\infty\}\). The frozen journal identifier is\par
https://doi.org/10.4171/JEMS/1631.\par
The reviewed PDF is arXiv version 3, submitted on 5 June 2024 at 17:20:07 UTC, with 22 pages. The first-page stamp agrees. The publisher records submission on 3 June 2023, acceptance on 18 April 2024, and online publication on 8 April 2025. The frozen 2025 membership and DOI are unchanged. This author manuscript is not presumed identical to the journal version of record. The publisher PDF route required a subscription. An unversioned 20-page PDF linked from the author’s site was observed but was not substituted for the frozen source.\par
All source locations below refer to the printed pages of this 22-page PDF. It has 335 387 bytes and SHA256\par
`456ea221e1b2e9cd4bfdbe0b2b825ec3eb9296be5503709e2b91055be12e8d51`.\par
The main result concerns the semisimplification of cohomology over a number field and functions on the proalgebraic completion. Full extensions and proreductive coefficients require separate arguments. The findings below distinguish false statements and proof steps from conclusions that remain valid after repair.\par
\section*{Confirmed problems}
\subsection*{P1. Ramified stalks lack the asserted simultaneous equivariance (Major)}
Location. Proposition 3.3, pp. 8–9, especially the stalk inclusion and its claimed \(H\)-stability on p. 9. This bridge is used in Corollary 3.4 and the main induction on p. 15.\par
Example under the actual hypotheses. Take \(F=\mathbb Q\) and \[X=\mathbb P^1\setminus\{1/3,\infty\},\quad x=1,\quad
 f(z)=\frac{27}{4}z(z-1)^2 .\] This is a smooth normal quasi-projective curve with rational point \(x\). The map is quasi-finite on \(X\), finite étale above \(\mathbb P^1\setminus\{0,1,\infty\}\), and sends \(x\) to zero. Let \(\mathcal L\) be the rank-one sign local system of the étale cover \(u^2=z-1/3\) of \(X\), and use the unit tangent at zero downstairs. For the quadratic character \(\chi_d\) of \(\mathbb Q(\sqrt d)/\mathbb Q\), the stalk \(\mathcal L_x\) is \(\chi_6\).\par
There are two ramified branches at \(z=1\) and one unramified branch at \(z=0\). Since \[f(1+h)=\frac{27}{4}h^2+\frac{27}{4}h^3,\] the ramified branches require \(\mathbb Q(\sqrt3)\). Their Galois module is \(\chi_6\otimes(1\oplus\chi_3)=\chi_6\oplus\chi_2\). The third branch contributes \(\chi_{-3}\). Thus \[(f_*\mathcal L)_{0_v}
     =\chi_6\oplus\chi_2\oplus\chi_{-3}.\] The characters are distinct, so the \(G_\mathbb Q\)-equivariant image of \(\mathcal L_x\) is unique. In branch coordinates it is the line spanned by \(e_1+e_2\), supported on both ramified branches.\par
The cubic has geometric permutation monodromy \(S_3\): it is connected of degree three and its discriminant \(-19683\,t(t-1)/16\) is not a square in \(\overline{\mathbb Q}(t)\). The subgroup \(H\) for the first branch maps onto its point stabilizer and hence contains an element inducing \((23)\). This sends the support \(\{1,2\}\) to \(\{1,3\}\). Possible Kummer signs do not change the support. The unique Galois-equivariant line is therefore not \(H\)-stable.\par
Counterevidence, impact, and boundary. A single branch gives an \(H\)-stable stalk that is not Galois-stable; the diagonal over the ramified branches is Galois-stable but not \(H\)-stable. The printed endomorphism-space argument needs both. This false intermediate assertion invalidates a central derivation, warranting Major severity. It does not disprove Proposition 3.3 or Theorem 1.2.\par
A replacement is available after the independent Artin argument in Section 3. Fix the target tangent \(v\) and use the two tangents \(v_1,v_2\) of Proposition 4.1. After a finite extension \(K/F\) a tangent above \(v_1\) exists. The valid finite-index argument gives \(W\twoheadrightarrow V|_{G_K}\) with \(W\) a finite-dimensional \(G_K\)-stable subspace of the function ring at \(v_1\). Its \(G_F\)-span \(W'\) is finite-dimensional, and \[\operatorname{Ind}_{G_K}^{G_F}W
       \subset W'\otimes\mathbb Q_p[G_F/G_K]
       \quad\hbox{has a quotient }V .\] The Artin factor occurs at \(v_2\); Proposition 4.1 then realizes \(V\) at \(v\). This uses neither a simultaneous stalk inclusion nor membership at every \(K\)-rational tangent. Finite-index induction of geometrically semisimple representations gives the same repair for the proreductive version.\par
\subsection*{P2. Full K3 cohomology is reducible (Major)}
Location. Corollary 9.7, p. 19, final paragraph; compare Proposition 9.6 and the cited Huybrechts Corollary 6.4.7.\par
In polarized K3 moduli the polarization class is a nonzero invariant line in the rank-22 system \(R^2\pi_*\mathbb Q_p(1)\). Its primitive orthogonal complement has rank 21. A smooth polarized quartic, for example, has the nonzero hyperplane class of square four. The full representation is therefore reducible. The cited monodromy result fixes this class and concerns the primitive part, as the public Huybrechts draft, PDF p. 118, explicitly states.\par
Proposition 9.6 cannot be applied to the full \(H^2\) as printed. Knowing the endomorphisms of the primitive part alone does not supply the off-diagonal terms in \(\operatorname{End}(H^2)\). This invalidates the stated proof of a substantive K3 application and is Major, although the conclusion has the following repair.\par
After the repaired proreductive bridge in P1, write \[H^2(X_{\overline F},\mathbb Q_p)(1)=\mathbb Q_p\oplus M,\qquad \dim M=21 .\] The primitive variation is geometrically absolutely irreducible, so \(\operatorname{End}(M)\in\mathcal C_F^{\mathrm{red}}\). The polarization identifies \(M^\vee\simeq M\). Pass to a finite extension killing the orthogonal determinant character. Then \[\Lambda^2M\ \text{is a subquotient of }\operatorname{End}(M),\qquad
 (\Lambda^2M)^{\otimes10}\twoheadrightarrow
 \Lambda^{20}M\simeq M^\vee\otimes\det M\simeq M .\] Tensor and subquotient closure give \(M\in\mathcal C_F^{\mathrm{red}}\). Hence \[\operatorname{End}(\mathbb Q_p\oplus M)=\mathbb Q_p\oplus M\oplus M\oplus\operatorname{End}(M)\] is in that class, and finite-extension descent finishes. This supplies the missing terms without presuming arbitrary dual or Tate closure of \(\mathcal C_F^{\mathrm{red}}\). The conclusion is not refuted.\par
\subsection*{P3. A level structure does not give a principal polarization (Major)}
Location. The abelian-variety paragraph of Corollary 9.7, p. 19. Choosing a basis of \(A[3]\) is said to give a point corresponding to \(A\) in the principally polarized moduli space.\par
An explicit obstruction is obtained from \(J=\operatorname{Jac}(y^2=x^5-x-1)\). The polynomial is irreducible modulo 3 and modulo 2 factors as \[(x^2+x+1)(x^3+x^2+1).\] Its Galois group is a transitive degree-five group containing an element of cycle type \((2)(3)\), hence a transposition. A primitive group containing a transposition is \(S_5\). Zarhin’s theorem for this characteristic-zero hyperelliptic Jacobian gives \(\operatorname{End}(J_{\overline{\mathbb Q}})=\mathbb Z\), and thus \(\operatorname{NS}(J_{\overline{\mathbb Q}})=\mathbb Z[\Theta]\).\par
Choose a subgroup \(H\) generated by a nonzero point of \(J[2]\), let \(q:J\to A=J/H\), and extend the number field so that \(H\) and a full level-three structure of \(A\) are defined. If \(A\) had a principal polarization \(\lambda\), its pullback would be \(n\Theta\) for a positive integer \(n\). On the surface \(J\) its Euler characteristic would be both \[\chi(q^*\lambda)=\deg(q)\chi(\lambda)=2
       \quad\hbox{and}\quad \chi(n\Theta)=n^2,\] a contradiction. No further extension changes this geometric obstruction. A level-three basis does not make this \(A\) principally polarized.\par
This missing reduction occurs at the realization step of the abelian-variety assertion in Corollary 9.7. The stated arbitrary abelian variety need not define a point of the principally polarized moduli space, even after a finite extension. Consequently, the displayed application of Proposition 9.6 does not establish that assertion as written. The defect is Major for this derivation.\par
Passing over a finite extension to a principally polarized abelian variety isogenous to \(A\), and choosing a symplectic level structure there, repairs the argument because rational \(H^1\) is invariant under isogeny. The corollary itself is not refuted.\par
\subsection*{P4. Nonnormal finite-index subgroups have no quotient group (Minor)}
Location. Remark 2.3, p. 5. Lemma 2.2 assumes only that \(\Gamma_1\) is open of finite index, but the remark calls \(\Gamma/\Gamma_1\) a constant finite group scheme. For the discrete finite group \(S_3\) and \(\Gamma_1=\langle(12)\rangle\), all hypotheses hold and the index is three, while \(\Gamma_1\) is not normal. It cannot be the kernel of a quotient group homomorphism. The homogeneous quotient is a constant finite scheme. Use scheme, or add normality for the group statement. The coordinate restriction lemma and its applications survive.\par
\subsection*{P5. Incorrect descent of the cyclic proalgebraic completion (Minor)}
Location. Example 2.4, pp. 5–6, the \(\mathbb Q_p\)-product \(\mathbb G_a\times H_0\times T_0\) and the module \(\widehat{\mathbb Z}(1)\) assigned to \(H_0\).\par
Every finite quotient of the completion of the discrete group \(\mathbb Z\) is generated by the rational image of \(1\); its points are therefore rational. The maximal profinite quotient is constant \(\widehat{\mathbb Z}\), not its cyclotomic twist. A quotient \(\mathbb Z/n\) with nontrivial local cyclotomic action modulo \(n\) distinguishes these group schemes.\par
Even after correcting the twist, the asserted rational product does not follow. The character sequence \[0\longrightarrow\mu_\infty\longrightarrow
 \overline{\mathbb Q}_p^\times\longrightarrow
 \overline{\mathbb Q}_p^\times/\mu_\infty\longrightarrow0\] has no \(G_{\mathbb Q_p}\)-equivariant splitting. The class \([p^{1/n}]\) is Galois-fixed for every \(n>1\), because its conjugates differ by roots of unity. An equivariant section would send it to an element of \(\mathbb Q_p^\times\) with normalized valuation \(1/n\), impossible. An abstract splitting over \(\overline{\mathbb Q}_p\) need not descend.\par
The safe rational description uses the additive factor and the multiplicative-type group with character module \(\overline{\mathbb Q}_p^\times\) and its actual Galois action. This illustrative description is not needed by the subsequent matrix-coefficient proof, which explains the limited severity.\par
\subsection*{P6. Finitely many roots of unity do not force the claimed local-finiteness criterion (Minor)}
Location. Example 2.6, p. 7. Take the allowed characteristic-zero field \(K=\mathbb R\), which has only two roots of unity, and \(X=\mathbb G_m\), \(x=1\). Its absolute Galois group has order two, so every function in the geometric proalgebraic coordinate ring has a finite-dimensional Galois orbit.\par
The continuous character \[\widehat{\mathbb Z}\twoheadrightarrow\mathbb Z_p
 \longrightarrow\mathbb Q_p^\times,\qquad
 t\longmapsto(1+p)^t\] has infinite image; for \(p=2\) use \(5^t\). It is a regular matrix coefficient. Complex conjugation interchanges this function with its inverse. It cannot factor through \(\mathbb G_a\times\widehat{\mathbb Z}\): a group-like function factoring there is an algebraic character, trivial on \(\mathbb G_a\) and of finite image on the profinite factor. This contradicts the infinite image.\par
The printed equivalence is therefore false under its stated field hypothesis. An infinite \(p\)-adic cyclotomic image is the relevant additional condition. Number fields satisfy it, so this example does not contradict the universality theorem.\par
\subsection*{P7. Faithfulness is insufficient for one high tensor power (Major)}
Location. Lemma 5.1, p. 12. The one-dimensional sign representation of \(C_2\) is faithful, but every tensor power is one-dimensional and cannot contain the two-dimensional regular representation. This meets the printed finite-group and characteristic-zero hypotheses.\par
The failure occurs in the actual proof: the representations \(W_i\) are only required to be faithful and of bounded dimension, and the printed construction does not require a trivial summand to be retained. The sign representation meets these conditions and can occur as a subrepresentation of the faithful permutation module. The invalid tensor-power assertion supplies the regular representation used to conclude Lemma 5.1. That lemma is used for finite-extension descent in Corollary 5.2 and explicitly as the dimension-zero base case of Theorem 1.2 on p. 15. This is a Major defect in the core derivation.\par
Retaining a trivial summand allows tensor degrees to be padded. The explicit permutation construction in Section 3 below gives a bounded power containing the regular image representation. These repairs do not refute the Artin conclusion. The earlier p. 12 inference from an action kernel contained in \(G_K\) to a quotient with kernel exactly \(G_K\) is not justified by containment alone; Section 3 bypasses that proof obligation. No counterexample to the geometric Artin conclusion is claimed.\par
\subsection*{P8. Arbitrary-field scope fails in characteristic \(p\) (Minor)}
Location. Proposition 7.1, p. 14, including the finite-dimensionality claim for continuous cocycles. The preceding setup on p. 6 allows any base field.\par
Take \(K=\mathbb F_p\), \(X=\mathbb A^1_K\) with base point zero, and \(\mathcal L=\mathbb Q_p\). Let \(k=\overline{\mathbb F}_p\), \(A=W(k)\), and \(E=A[1/p]\), with Witt Frobenius \(\sigma\). Artin–Schreier–Witt theory gives an additive, Frobenius-equivariant map \[A\longrightarrow H^1_{\mathrm{cont}}(\pi_1(X_k,0),\mathbb Z_p),
           \qquad a\longmapsto[a[x]].\] Indeed the compatible finite-level classes lie in \(W_n(k[x])/(F-1)\) and define continuous characters. For a unit \(a\), the reduction is \(a_0x\), which is not \(b^p-b\): a nonconstant polynomial \(b\) gives degree divisible by \(p\), whereas \(a_0x\) has degree one. For any nonzero \(a\), divide by its finite power of \(p\) and use torsion-freeness of \(\operatorname{Hom}_{\mathrm{cont}}(\pi_1,\mathbb Z_p)\). The map is therefore injective, and rationalization embeds \(E\) into \(H^1(X_k,\mathbb Q_p)\).\par
This subspace is not locally finite for Frobenius. For \(n\ge1\) set \[D_n(a)=\det\bigl(\sigma^{i+j}a\bigr)_{0\le i,j<n}.\] An integral normal-basis element \(b\) of the degree-\(n\) unramified extension of \(\mathbb Q_p\) has \(D_n(b)\ne0\). Thus \(D_n(a+tb)\), as a polynomial in \(t\in\mathbb Q_p\), has nonzero leading coefficient. The set \(D_n\ne0\) is open and dense in the complete metric ring \(A\). Baire’s theorem supplies an \(a\) for which every \(D_n(a)\) is nonzero. Any finite \(\mathbb Q_p\)-linear relation among its Frobenius iterates would force one of these determinants to vanish. Its orbit therefore spans an infinite-dimensional space. The inverse-Frobenius convention gives the same conclusion.\par
The right-hand side of Proposition 7.1, with this trivial stalk, is locally finite by definition, and all its subquotients are locally finite. It cannot have the displayed \(H^1\) as a subquotient. This is a counterexample to the proposition under all its written hypotheses, not merely to finite-dimensionality in the proof. Require the needed finiteness and local-finiteness assumptions, in particular the characteristic-zero setting of Section 8, or an appropriate invertible-coefficient setting. The actual number-field uses avoid the counterexample; the defect is in the added arbitrary-field generality.\par
\subsection*{P9. Removing one fibre need not remove ramification (Minor)}
Location. Lemma 10.1(ii), p. 20. Over \(\mathbb Q\) take \(D=C=\mathbb P^1\), \(f(t)=t^2\), \(x=1\), and a nonzero tangent at \(x\). The map is finite surjective and even unramified at \(x\). Nevertheless \[D\setminus f^{-1}(f(x))=\mathbb P^1\setminus\{-1,1\}\simeq\mathbb G_m,
 \qquad C\setminus\{f(x)\}\simeq\mathbb A^1.\] The asserted map on geometric fundamental groups is \(\widehat{\mathbb Z}\to1\), which is not injective; the arithmetic map cannot be injective either. Require the restricted map to be finite étale, or delete all branch values and their preimages. The Belyi applications already delete \(\{0,1,\infty\}\), so their finite-index step is unaffected. Existence of a compatible tangent is a local fact and does not establish the claimed global open-subgroup property.\par
\section*{A noncircular Artin replacement}
This argument supplies the independent input needed in P1, without invoking Proposition 3.3 or Corollary 3.4.\par
Fix a finite Galois extension \(K/F\) with group \(G\). Choose an elliptic curve over \(K\) whose \(j\)-invariant generates \(K\) over \(\mathbb Q\), and a Belyi cover over \(K\). Let \(H\) be its open subgroup in the geometric free profinite group \(\Gamma\) of rank two. Preserving the conjugacy class of \(H\) gives an isomorphism of the proper curves and fixes that \(j\)-invariant; hence its stabilizer in \(G_F\) is \(G_K\).\par
There are finitely many open subgroups of each bounded index in \(\Gamma\). Intersect those of index at most \([\Gamma:H]\) to obtain a characteristic open subgroup \(N\subset H\). At each tangent \(0_w\) put \(S_w=\Gamma_w/N_w\). These have the same abstract finite group \(S\), with different possible Galois actions. The image \(A_w\subset\operatorname{Aut}(S_w)\) maps onto \(G\): an element in the action kernel fixes \(H\) and hence belongs to \(G_K\). The permutation module \(W_w=\mathbb Q[S_w]\) is a subspace of locally constant functions at \(0_w\) and has a trivial line.\par
The explicit construction in Lemma 4.2 permits child tangents \[w_1=\frac4{27}w,\qquad w_2=\frac13w .\] Here \(f'(0)=27/4\) and \(h(z)=-3z/(z-4/3)\) maps \((0,1/3,4/3)\) to \((0,1,\infty)\) with \(h'(0)=9/4\). At depth \(n\) in the iterated binary construction, the tangent \((4/27)^i(1/3)^{n-i}w\) occurs \(\binom ni\) times.\par
The tensor construction can be made an actual injection on locally constant rational functions. Pullback along the surjective product map is injective, and extension by zero from the \(G_F\)-stable open subgroup is an equivariant splitting of restriction. Thus the tensor of all leaf modules embeds in the locally constant functions at the root.\par
Choose a repeated leaf with multiplicity \(m\ge |S|\). In \(W_w^{\otimes m}=\mathbb Q[S_w^m]\), a tuple whose first \(|S|\) entries enumerate \(S_w\) has trivial \(A_w\)-stabilizer. Its orbit spans the regular representation of \(A_w\), which contains the inflated regular representation of \(G\). Arbitrary tails give \(|S|^{m-|S|}\) distinct free orbits. Taking invariant constant lines in all other leaf factors therefore produces arbitrarily many regular copies at the root as \(m\) increases. Every finite-image rational representation embeds in a finite sum of regular copies. This proves the stronger locally constant embedding claimed in Proposition 1.1, as well as Lemma 5.1 after extending coefficients. Locally constant functions lie in the proreductive ring too. Only Belyi’s theorem, finite generation, and the tensor part of Lemma 4.2/Proposition 4.1 enter; there is no circular use of P1.\par
\section*{Further proof audit}
\subsection*{Finite translates, finite index, and direct sums}
Lemmas 2.1–2.2 and 2.7–2.9 were checked using representative functions, finite-dimensional translation orbits, and finite-index induction. The span \(\mathcal F(\mathcal L)\) of monodromy operators must be distinguished from its dual space of functions.\par
For direct sums in Proposition 4.1, the quotients \(\Gamma_w\twoheadrightarrow\mathbb Z/2^m(1)\) have \(m+1\) distinct additive orders. Their order indicators are independent Galois-fixed locally constant functions. Hence there are arbitrarily large trivial subrepresentations at a fixed tangent. Given a finite subrepresentation \(A\) yielding \(V\), choose a fixed line disjoint from \(A\); \(A\oplus\mathbb Q_p\) yields \(V\oplus\mathbb Q_p\). This supplies the step without relying merely on large trivial subquotients.\par
\subsection*{The Tate character and semisimple induction}
In the version-3 proof of Lemma 6.2, nonzero unipotent Legendre monodromy gives an invariant line \(\chi_1\) and quotient \(\chi_1(-1)\). The operators from quotient to invariant line form \(\mathbb Q_p(1)\subset\operatorname{End}(\mathcal L_{0_v})\). Geometric absolute irreducibility embeds these coefficients in the function ring. The remaining CM/abelian-character classification is the cited Fontaine–Mazur input, not an independently reproved theorem here.\par
On p. 15 the induction supplies \(B^{\mathrm{ss}}\) for \(B=(R^j\pi_*\mathbb Q_p)_x\), but the last paragraph uses \(B\) itself and concludes membership of the full \(H^1\). No counterexample to this stronger membership is claimed. The stated semisimple goal follows by choosing a finite-dimensional coefficient submodule \(A_0\) supporting the finite-dimensional quotient in Proposition 7.1, and filtering \(A_0\), \(B\), and \(A_0\otimes B\) by composition series. Their constituents occur in tensor products of constituents. Subquotient, tensor, and direct-sum closure give the semisimplification of \(H^1\), precisely what Leray needs.\par
With this formulation, the reductions through cohomological descent, smooth schemes, Gysin/purity, finite étale trace splittings, and an alteration to a proper smooth fibration over an open of \(\mathbb P^1\) were checked. Artin vanishing on the affine base leaves \(H^0\) and \(H^1\). The underlying descent, alteration, and purity theorems remain identified foundational dependencies; their complete source literature has not been independently reproved.\par
\subsection*{Weights and proreductive coefficients}
The apparent sign discrepancy in Proposition 9.1 is not an error: nonpositive weights on \(\mathcal F(\mathcal L)\) become nonnegative weights on its dual function space. The \(\overline{\mathbb Q}_p\) coefficient field in Lemma 9.5 and absolute irreducibility in the \(\mathbb Q_p\) clause of Proposition 9.6 were checked on actual images. Text extraction must not erase them. Corollary 9.7 itself is stated with \(\overline{\mathbb Q}_p\) coefficients. The polarization arguments in P2–P3 were written over \(\mathbb Q_p\) and remain valid after that coefficient extension.\par
The companion Petrov manuscript arXiv:2012.13372v3, PDF pp. 17–18, includes tangential base points in the de Rham and almost-everywhere-unramified input. Its exact retrieved version is recorded separately, without equating it to the later Duke version. Lemma 9.3 is a conditional Fontaine–Mazur equivalence, not a proof of that conjecture. The full Deligne, Lafforgue, and Fontaine–Mazur source proofs were not independently reconstructed in this review.\par
In Lemma 10.1(i), an étale cover extending over \(F[[t]]\) is constant, not necessarily a trivial disjoint union over \(F\): \(\mathbb Q(\sqrt2)[[t]]\) illustrates this. The geometric-fibre comparison used there is correct, so this wording is not counted as another scientific defect.\par
\section*{Supporting checks}
The initial root reviewer read all 1 369 normalized source lines and all 452 manifest lines in contiguous ranges ending with zero remaining. Rendering alone was not inspection. The additional PDFs were fully parsed and hashed; reading of the 449-page Huybrechts draft was limited to the relevant pages, not a scientific audit of the whole book.\par
Three native auxiliaries supplied bounded checks during the initial review. `/root/check\_0251\_lemma23` covered Sections 2–4 and the ramified-stalk example; `/root/check\_0251\_section5` covered Sections 5–8, the Artin replacement, and the characteristic-\(p\) example; `/root/check\_0251\_theorem414` covered Sections 9–10 and the cited external dependencies. The initial root integrated that report. The current focused independent acceptance has separately attributed records and native read-only contributions in the verification materials. Neither stage is represented as human review.\par
The Personal-account observation is attributed to the external controller/user, not independently observed here. The current conversation URL was unavailable; no acquisition conversation URL was substituted.\par
Bounded arXiv, publisher, and author-page searches located no dedicated supplement or author code. This is a search limit, not a defect. No author code was executed. Own exact calculations and their raw failed serialization attempt were preserved beside the successful corrected output. The TeX drafting failure before file creation is separately recorded. The current report is compiled with the previously authorized portable Tectonic route; current compilation and actual page-inspection receipts are separate.\par
Assessment: four Major and five Minor problems. The universality theorem and K3 conclusion are not refuted; explicit repairs and their dependencies are supplied. Scientific decision: pending. Human reviews: empty. The current focused independent acceptance is recorded separately from the initial review and does not constitute final adjudication.\par
9 A. Petrov, Universality of the Galois action on the fundamental group of \(\mathbb P^1\setminus\{0,1,\infty\}\), arXiv:2109.09301v3, 5 June 2024. https://arxiv.org/abs/2109.09301v3. EMS Press, article identity and publication history, DOI 10.4171/JEMS/1631. https://ems.press/journals/jems/articles/14298702. A. Petrov, Geometrically irreducible \(p\)-adic local systems are de Rham up to a twist, arXiv:2012.13372v3, exact companion manuscript, PDF pp. 17–18. https://arxiv.org/abs/2012.13372v3. D. Huybrechts, Lectures on K3 Surfaces, public author draft, exact retrieved 449-page PDF, Corollary 6.4.7, PDF p. 118. Source URL, hash, and page-view record are retained. Yu. G. Zarhin, Hyperelliptic Jacobians without complex multiplication, arXiv:math/9909052v2, Theorem 2.1. The exact retrieved manuscript and relevant pages are retained. https://arxiv.org/abs/math/9909052v2.\par
\end{document}
